\documentclass[11pt,a4paper]{article}
\usepackage[utf8]{inputenc}
\usepackage[T1]{fontenc}
\usepackage[catalan]{babel}
\usepackage{lmodern}
\usepackage[sfdefault,lining]{FiraSans}
\usepackage[a4paper,top=2.3cm,bottom=2.2cm,left=1.8cm,right=1.8cm,headheight=1.3cm,headsep=0.4cm,footskip=1cm]{geometry}
\usepackage{amsmath,amssymb}
\usepackage{icomma}
\usepackage{graphicx}
\usepackage[table]{xcolor}
\usepackage{tikz}
\usepackage{fontawesome5}
\usepackage{ccicons}
\usepackage{fancyhdr}
\usepackage{lastpage}
\usepackage{tasks}
\usepackage{needspace}
\usepackage[most]{tcolorbox}
\usepackage{array}
\usepackage[hidelinks,pdftitle={Solucions de la fitxa de la unitat 7: Geometria bàsica del pla. 1r d'ESO},pdfsubject={Matemàtiques 1r ESO},pdfkeywords={punt, recta, semirecta, segment, rectes perpendiculars, mediatriu, angles, bisectriu, sistema sexagesimal, complementaris, suplementaris, polígons, diagonals, triangles, quadrilàters, polígons regulars}]{hyperref}

\definecolor{navy}{HTML}{1B2748}
\definecolor{accent}{HTML}{E8603A}
\definecolor{teal}{HTML}{1F9E90}
\definecolor{soft}{HTML}{EEF1F8}
\definecolor{muted}{HTML}{5D6685}

\setlength{\parindent}{0pt}
\setlength{\parskip}{3pt}
\settasks{label=\alph*), label-format=\bfseries\color{navy}, label-width=1.3em, item-indent=1.8em, column-sep=0.6em, after-item-skip=5pt, before-skip=3pt, after-skip=4pt}

% ---------- capçalera i peu ----------
\newcommand{\llicencia}{\href{https://creativecommons.org/licenses/by-nc-sa/4.0/deed.ca}{\ccbyncsa}\ \ Llicència Creative Commons Reconeixement-NoComercial-CompartirIgual 4.0 Internacional (CC BY-NC-SA 4.0)}
\pagestyle{fancy}
\fancyhf{}
\fancyhead[L]{\raisebox{-0.35\height}{\includegraphics[height=0.95cm]{logo-gva.png}}}
\fancyhead[C]{\small\color{muted}Solucions de la fitxa · Unitat 7 · 1r d'ESO}
\fancyhead[R]{\raisebox{-0.35\height}{\includegraphics[height=0.95cm]{logo-vermellar.png}}}
\fancyfoot[L]{\scriptsize\color{muted}\llicencia}
\fancyfoot[R]{\small\color{muted}\thepage}
\renewcommand{\headrulewidth}{0.4pt}
\renewcommand{\headrule}{\color{muted!40}\hrule height \headrulewidth}
\fancypagestyle{primera}{\fancyhf{}\renewcommand{\headrulewidth}{0pt}
  \fancyfoot[L]{\scriptsize\color{muted}\llicencia}
  \fancyfoot[R]{\small\color{muted}\thepage}}

% ---------- nivells i exercicis ----------
\newcommand{\estrelles}[1]{{\color{accent}\ifcase#1\or\faStar\faStar[regular]\faStar[regular]\or\faStar\faStar\faStar[regular]\or\faStar\faStar\faStar\fi}}
\newcounter{exer}
\newcommand{\ex}[1]{\par\needspace{4\baselineskip}\medskip\stepcounter{exer}\noindent{\color{navy}\bfseries\theexer.}\ \ {\small\estrelles{#1}}\ \ }

\newcommand{\seccio}[2]{\par\needspace{10\baselineskip}\bigskip
  \begin{tcolorbox}[colback=navy,colframe=navy,coltext=white,arc=2mm,boxrule=0pt,left=3mm,right=3mm,top=1.2mm,bottom=1.2mm]
  {\large\bfseries #1\quad #2}
  \end{tcolorbox}\vspace{-2pt}}

\newcolumntype{C}[1]{>{\centering\arraybackslash}p{#1}}

% quadre de recordatori (només a la fitxa d'exercicis)
\newcommand{\recorda}[1]{\par\smallskip\begin{tcolorbox}[colback=accent!7,colframe=accent!55,arc=2mm,boxrule=0.5pt,left=2.5mm,right=2.5mm,top=1mm,bottom=1mm]\small{\color{accent}\faLightbulb[regular]}\ \textbf{Recorda:}\ #1\end{tcolorbox}}
\newcommand{\res}[1]{\boxed{#1}}\newcommand{\rest}[1]{\fbox{#1}}

\def\fmneed{14}
% =====================================================================
% Capçalera FentMates de les fitxes: un color per a cada curs.
% Els colors són els de la pàgina inicial (LEVELS a assets/js/materials.js);
% el logo ve de _eines/logo/pdf/ (python3 _eines/logo/genera.py).
%
% Ús, al final del preàmbul de la unitat:
%   % =====================================================================
% Capçalera FentMates de les fitxes: un color per a cada curs.
% Els colors són els de la pàgina inicial (LEVELS a assets/js/materials.js);
% el logo ve de _eines/logo/pdf/ (python3 _eines/logo/genera.py).
%
% Ús, al final del preàmbul de la unitat:
%   % =====================================================================
% Capçalera FentMates de les fitxes: un color per a cada curs.
% Els colors són els de la pàgina inicial (LEVELS a assets/js/materials.js);
% el logo ve de _eines/logo/pdf/ (python3 _eines/logo/genera.py).
%
% Ús, al final del preàmbul de la unitat:
%   \input{../comu/fentmates.tex}
%   \fmcurs{1eso}
% i, al començament del document:
%   \capcalera{Matemàtiques}{1r ESO}{Fitxa d'exercicis · Unitat 2}{Divisibilitat}
% Opcional abans de \input: \def\fmneed{14} (espai mínim abans de cada secció).
% =====================================================================
\makeatletter
\def\fm@logodir{../../logo/pdf}
\providecommand{\fmneed}{10}
\def\fm@cursnom{}

% clau del curs, color i color del text damunt del color
\newcommand{\fm@curs}[3]{\@namedef{fm@color@#1}{#2}\@namedef{fm@text@#1}{#3}}
\fm@curs{1eso}{E6B84A}{navy}
\fm@curs{2eso}{3B82F6}{white}
\fm@curs{3eso}{FF7A50}{white}
\fm@curs{4eso-a}{A78BFA}{white}
\fm@curs{4eso-b}{6D28D9}{white}
\fm@curs{1batx-mi}{9B1C1C}{white}
\fm@curs{1batx-ccss}{EF6B6B}{white}
\fm@curs{2batx-mii}{1E6B47}{white}
\fm@curs{2batx-ccss}{2BB3A3}{white}

\newcommand{\fmcurs}[1]{%
  \@ifundefined{fm@color@#1}{\PackageError{fentmates}{Curs desconegut: #1}{}}{}%
  \edef\fm@tmp{\noexpand\definecolor{curs}{HTML}{\@nameuse{fm@color@#1}}%
               \noexpand\colorlet{curstext}{\@nameuse{fm@text@#1}}}\fm@tmp
  \colorlet{cursfosc}{curs!72!black}%
  \def\fm@logo{\fm@logodir/vertical-#1.pdf}%
  \def\fm@marca{\fm@logodir/marca-#1.pdf}}

% ---------- capçalera de la primera pàgina ----------
% #1 matèria (petit), #2 curs, #3 tipus de document i unitat, #4 títol
\newcommand{\capcalera}[4]{%
  \gdef\fm@cursnom{#2}%
  \par\noindent
  \begin{tcolorbox}[colback=curs,colframe=curs,boxrule=0pt,arc=2.2mm,width=3.5cm,height=2.15cm,
      valign=center,left=3mm,right=2mm,top=1mm,bottom=1mm,nobeforeafter,box align=center]
    \color{curstext}{\footnotesize #1}\par\vspace{1pt}{\fontsize{19}{21}\selectfont\bfseries #2\par}%
  \end{tcolorbox}\hspace{5mm}%
  \begin{minipage}[c]{\dimexpr\linewidth-3.5cm-5mm-2.9cm\relax}
    {\color{cursfosc}\small\bfseries\MakeUppercase{#3}}\par\vspace{2pt}
    {\color{navy}\fontsize{20}{23}\selectfont\bfseries #4\par}
  \end{minipage}\hfill
  \begin{minipage}[c]{2.6cm}\raggedleft\includegraphics[width=2.6cm]{\fm@logo}\end{minipage}\par
  \vspace{2pt}}

% ---------- barres de secció del color del curs ----------
\renewcommand{\seccio}[2]{\par\needspace{\fmneed\baselineskip}\bigskip
  \begin{tcolorbox}[colback=curs,colframe=curs,coltext=curstext,arc=2mm,boxrule=0pt,left=3mm,right=3mm,top=1.2mm,bottom=1.2mm]
  {\large\bfseries #1\quad #2}
  \end{tcolorbox}\vspace{-2pt}}

% ---------- estrelles de dificultat del color del curs (a Batxillerat no se'n posen) ----------
\@ifundefined{estrelles}{}{%
  \renewcommand{\estrelles}[1]{{\color{curs}\ifcase#1\or\faStar\faStar[regular]\faStar[regular]\or\faStar\faStar\faStar[regular]\or\faStar\faStar\faStar\fi}}}

% ---------- capçalera i peu de les pàgines ----------
% primera pàgina: logos institucionals al peu
\fancypagestyle{primera}{\fancyhf{}\renewcommand{\headrulewidth}{0pt}
  \fancyfoot[L]{\raisebox{-0.32\height}{\includegraphics[height=0.78cm]{logo-gva.png}}\hspace{5mm}%
    \raisebox{-0.32\height}{\includegraphics[height=0.78cm]{logo-vermellar.png}}}
  \fancyfoot[R]{\scriptsize\color{muted}\href{https://creativecommons.org/licenses/by-nc-sa/4.0/deed.ca}{\ccbyncsa}\ \ CC BY-NC-SA 4.0 · fentmates.org\qquad\small\thepage}}
% la resta: el símbol a l'esquerra i el curs a la dreta, amb el títol al centre com abans
\fancyhead[L]{\raisebox{-0.3\height}{\includegraphics[height=0.62cm]{\fm@marca}}}
\fancyhead[R]{\small\bfseries\color{cursfosc}\fm@cursnom}
\renewcommand{\headrule}{\color{curs!45}\hrule height \headrulewidth}
\makeatother

%   \fmcurs{1eso}
% i, al començament del document:
%   \capcalera{Matemàtiques}{1r ESO}{Fitxa d'exercicis · Unitat 2}{Divisibilitat}
% Opcional abans de \input: \def\fmneed{14} (espai mínim abans de cada secció).
% =====================================================================
\makeatletter
\def\fm@logodir{../../logo/pdf}
\providecommand{\fmneed}{10}
\def\fm@cursnom{}

% clau del curs, color i color del text damunt del color
\newcommand{\fm@curs}[3]{\@namedef{fm@color@#1}{#2}\@namedef{fm@text@#1}{#3}}
\fm@curs{1eso}{E6B84A}{navy}
\fm@curs{2eso}{3B82F6}{white}
\fm@curs{3eso}{FF7A50}{white}
\fm@curs{4eso-a}{A78BFA}{white}
\fm@curs{4eso-b}{6D28D9}{white}
\fm@curs{1batx-mi}{9B1C1C}{white}
\fm@curs{1batx-ccss}{EF6B6B}{white}
\fm@curs{2batx-mii}{1E6B47}{white}
\fm@curs{2batx-ccss}{2BB3A3}{white}

\newcommand{\fmcurs}[1]{%
  \@ifundefined{fm@color@#1}{\PackageError{fentmates}{Curs desconegut: #1}{}}{}%
  \edef\fm@tmp{\noexpand\definecolor{curs}{HTML}{\@nameuse{fm@color@#1}}%
               \noexpand\colorlet{curstext}{\@nameuse{fm@text@#1}}}\fm@tmp
  \colorlet{cursfosc}{curs!72!black}%
  \def\fm@logo{\fm@logodir/vertical-#1.pdf}%
  \def\fm@marca{\fm@logodir/marca-#1.pdf}}

% ---------- capçalera de la primera pàgina ----------
% #1 matèria (petit), #2 curs, #3 tipus de document i unitat, #4 títol
\newcommand{\capcalera}[4]{%
  \gdef\fm@cursnom{#2}%
  \par\noindent
  \begin{tcolorbox}[colback=curs,colframe=curs,boxrule=0pt,arc=2.2mm,width=3.5cm,height=2.15cm,
      valign=center,left=3mm,right=2mm,top=1mm,bottom=1mm,nobeforeafter,box align=center]
    \color{curstext}{\footnotesize #1}\par\vspace{1pt}{\fontsize{19}{21}\selectfont\bfseries #2\par}%
  \end{tcolorbox}\hspace{5mm}%
  \begin{minipage}[c]{\dimexpr\linewidth-3.5cm-5mm-2.9cm\relax}
    {\color{cursfosc}\small\bfseries\MakeUppercase{#3}}\par\vspace{2pt}
    {\color{navy}\fontsize{20}{23}\selectfont\bfseries #4\par}
  \end{minipage}\hfill
  \begin{minipage}[c]{2.6cm}\raggedleft\includegraphics[width=2.6cm]{\fm@logo}\end{minipage}\par
  \vspace{2pt}}

% ---------- barres de secció del color del curs ----------
\renewcommand{\seccio}[2]{\par\needspace{\fmneed\baselineskip}\bigskip
  \begin{tcolorbox}[colback=curs,colframe=curs,coltext=curstext,arc=2mm,boxrule=0pt,left=3mm,right=3mm,top=1.2mm,bottom=1.2mm]
  {\large\bfseries #1\quad #2}
  \end{tcolorbox}\vspace{-2pt}}

% ---------- estrelles de dificultat del color del curs (a Batxillerat no se'n posen) ----------
\@ifundefined{estrelles}{}{%
  \renewcommand{\estrelles}[1]{{\color{curs}\ifcase#1\or\faStar\faStar[regular]\faStar[regular]\or\faStar\faStar\faStar[regular]\or\faStar\faStar\faStar\fi}}}

% ---------- capçalera i peu de les pàgines ----------
% primera pàgina: logos institucionals al peu
\fancypagestyle{primera}{\fancyhf{}\renewcommand{\headrulewidth}{0pt}
  \fancyfoot[L]{\raisebox{-0.32\height}{\includegraphics[height=0.78cm]{logo-gva.png}}\hspace{5mm}%
    \raisebox{-0.32\height}{\includegraphics[height=0.78cm]{logo-vermellar.png}}}
  \fancyfoot[R]{\scriptsize\color{muted}\href{https://creativecommons.org/licenses/by-nc-sa/4.0/deed.ca}{\ccbyncsa}\ \ CC BY-NC-SA 4.0 · fentmates.org\qquad\small\thepage}}
% la resta: el símbol a l'esquerra i el curs a la dreta, amb el títol al centre com abans
\fancyhead[L]{\raisebox{-0.3\height}{\includegraphics[height=0.62cm]{\fm@marca}}}
\fancyhead[R]{\small\bfseries\color{cursfosc}\fm@cursnom}
\renewcommand{\headrule}{\color{curs!45}\hrule height \headrulewidth}
\makeatother

%   \fmcurs{1eso}
% i, al començament del document:
%   \capcalera{Matemàtiques}{1r ESO}{Fitxa d'exercicis · Unitat 2}{Divisibilitat}
% Opcional abans de \input: \def\fmneed{14} (espai mínim abans de cada secció).
% =====================================================================
\makeatletter
\def\fm@logodir{../../logo/pdf}
\providecommand{\fmneed}{10}
\def\fm@cursnom{}

% clau del curs, color i color del text damunt del color
\newcommand{\fm@curs}[3]{\@namedef{fm@color@#1}{#2}\@namedef{fm@text@#1}{#3}}
\fm@curs{1eso}{E6B84A}{navy}
\fm@curs{2eso}{3B82F6}{white}
\fm@curs{3eso}{FF7A50}{white}
\fm@curs{4eso-a}{A78BFA}{white}
\fm@curs{4eso-b}{6D28D9}{white}
\fm@curs{1batx-mi}{9B1C1C}{white}
\fm@curs{1batx-ccss}{EF6B6B}{white}
\fm@curs{2batx-mii}{1E6B47}{white}
\fm@curs{2batx-ccss}{2BB3A3}{white}

\newcommand{\fmcurs}[1]{%
  \@ifundefined{fm@color@#1}{\PackageError{fentmates}{Curs desconegut: #1}{}}{}%
  \edef\fm@tmp{\noexpand\definecolor{curs}{HTML}{\@nameuse{fm@color@#1}}%
               \noexpand\colorlet{curstext}{\@nameuse{fm@text@#1}}}\fm@tmp
  \colorlet{cursfosc}{curs!72!black}%
  \def\fm@logo{\fm@logodir/vertical-#1.pdf}%
  \def\fm@marca{\fm@logodir/marca-#1.pdf}}

% ---------- capçalera de la primera pàgina ----------
% #1 matèria (petit), #2 curs, #3 tipus de document i unitat, #4 títol
\newcommand{\capcalera}[4]{%
  \gdef\fm@cursnom{#2}%
  \par\noindent
  \begin{tcolorbox}[colback=curs,colframe=curs,boxrule=0pt,arc=2.2mm,width=3.5cm,height=2.15cm,
      valign=center,left=3mm,right=2mm,top=1mm,bottom=1mm,nobeforeafter,box align=center]
    \color{curstext}{\footnotesize #1}\par\vspace{1pt}{\fontsize{19}{21}\selectfont\bfseries #2\par}%
  \end{tcolorbox}\hspace{5mm}%
  \begin{minipage}[c]{\dimexpr\linewidth-3.5cm-5mm-2.9cm\relax}
    {\color{cursfosc}\small\bfseries\MakeUppercase{#3}}\par\vspace{2pt}
    {\color{navy}\fontsize{20}{23}\selectfont\bfseries #4\par}
  \end{minipage}\hfill
  \begin{minipage}[c]{2.6cm}\raggedleft\includegraphics[width=2.6cm]{\fm@logo}\end{minipage}\par
  \vspace{2pt}}

% ---------- barres de secció del color del curs ----------
\renewcommand{\seccio}[2]{\par\needspace{\fmneed\baselineskip}\bigskip
  \begin{tcolorbox}[colback=curs,colframe=curs,coltext=curstext,arc=2mm,boxrule=0pt,left=3mm,right=3mm,top=1.2mm,bottom=1.2mm]
  {\large\bfseries #1\quad #2}
  \end{tcolorbox}\vspace{-2pt}}

% ---------- estrelles de dificultat del color del curs (a Batxillerat no se'n posen) ----------
\@ifundefined{estrelles}{}{%
  \renewcommand{\estrelles}[1]{{\color{curs}\ifcase#1\or\faStar\faStar[regular]\faStar[regular]\or\faStar\faStar\faStar[regular]\or\faStar\faStar\faStar\fi}}}

% ---------- capçalera i peu de les pàgines ----------
% primera pàgina: logos institucionals al peu
\fancypagestyle{primera}{\fancyhf{}\renewcommand{\headrulewidth}{0pt}
  \fancyfoot[L]{\raisebox{-0.32\height}{\includegraphics[height=0.78cm]{logo-gva.png}}\hspace{5mm}%
    \raisebox{-0.32\height}{\includegraphics[height=0.78cm]{logo-vermellar.png}}}
  \fancyfoot[R]{\scriptsize\color{muted}\href{https://creativecommons.org/licenses/by-nc-sa/4.0/deed.ca}{\ccbyncsa}\ \ CC BY-NC-SA 4.0 · fentmates.org\qquad\small\thepage}}
% la resta: el símbol a l'esquerra i el curs a la dreta, amb el títol al centre com abans
\fancyhead[L]{\raisebox{-0.3\height}{\includegraphics[height=0.62cm]{\fm@marca}}}
\fancyhead[R]{\small\bfseries\color{cursfosc}\fm@cursnom}
\renewcommand{\headrule}{\color{curs!45}\hrule height \headrulewidth}
\makeatother

\fmcurs{1eso}

\begin{document}
\thispagestyle{primera}

\capcalera{Matemàtiques}{1r ESO}{Solucions de la fitxa · Unitat 7}{Geometria bàsica del pla}
\vspace{8pt}

\begin{tcolorbox}[colback=soft,colframe=soft,arc=2mm,left=3mm,right=3mm,top=1.5mm,bottom=1.5mm]
\small Intenta resoldre cada exercici abans de mirar-ne la solució. Si no t'ix el mateix resultat, repassa els passos o consulta els apunts interactius de la unitat: hi trobaràs una calculadora d'angles en el sistema sexagesimal i eines per a explorar rectes, angles, triangles i polígons.
\end{tcolorbox}

% =====================================================
\seccio{1}{Punts, rectes, semirectes i segments}

\ex{1} \begin{tasks}(1)
  \task $A$, $B$ i $C$ són de la recta $r$ (estan \rest{alineats}); $D$ \rest{no} hi és.
  \task \rest{$\overline{AB}$, $\overline{AC}$ i $\overline{BC}$} (3 segments).
  \task \rest{Dues}: una cap a $A$ i una altra cap a $C$.
\end{tasks}

\ex{1} \begin{tasks}(3)
  \task \rest{Recta} (dues fletxes)
  \task \rest{Segment} (dos extrems)
  \task \rest{Semirecta} (una fletxa)
\end{tasks}

\ex{1} \begin{tasks}(1)
  \task $7:2=\res{3,5\text{ cm}}$
  \task $2\cdot4,2=\res{8,4\text{ cm}}$
\end{tasks}

\ex{2} Amb 4 punts: $3+2+1=\res{6}$ segments. Amb 6 punts: $5+4+3+2+1=\res{15}$ segments.

\ex{2} En el primer cas, $\overline{AC}=5+3=\res{8\text{ cm}}$. En el segon, $\overline{AC}=5-3=\res{2\text{ cm}}$.

\ex{3} De cada punt ixen tantes rectes com punts queden, i cada recta es compta dues vegades: \ $\frac{4\cdot3}{2}=\res{6}$, \ $\frac{5\cdot4}{2}=\res{10}$, \ $\frac{7\cdot6}{2}=\res{21}$.


% =====================================================
\seccio{2}{Posicions relatives de rectes. Mediatriu}

\ex{1} \begin{tasks}(4)
  \task \rest{paral·leles}
  \task \rest{perpendiculars}
  \task \rest{perpendiculars}
  \task \rest{secants}
\end{tasks}

\ex{1} \begin{tasks}(2)
  \task \rest{Paral·leles}
  \task \rest{Perpendiculars}
  \task \rest{Perpendiculars} (angle de $90^\circ$)
  \task \rest{Secants} (angle de $60^\circ$)
\end{tasks}

\ex{1} Amb el compàs obert més de 3\,cm, traça arcs des d'$A$ i des de $B$ que es tallen per damunt i per davall del segment; la recta que unix els dos punts de tall és la mediatriu. Talla el segment en el punt mitjà, a \rest{3\,cm} d'$A$, i és perpendicular al segment.

\ex{2} \begin{tasks}(1)
  \task \rest{Cert}: es tallen en un punt.
  \task \rest{Fals}: només si formen angles rectes.
  \task \rest{Cert}.
  \task \rest{Fals}: $r$ i $t$ són paral·leles.
\end{tasks}

\ex{2} Els punts de la mediatriu equidisten dels extrems: $\overline{PB}=\res{5\text{ cm}}$. \ $\overline{AM}=8:2=\res{4\text{ cm}}$.

\ex{3} Els punts que equidisten de $A$ i $B$ són els de la \rest{mediatriu de $\overline{AB}$}. La font ha d'estar en el \rest{punt on la mediatriu talla el camí}.


% =====================================================
\seccio{3}{Angles: elements, mesura i classificació}

\ex{1} \begin{tasks}(4)
  \task \rest{agut}
  \task \rest{recte}
  \task \rest{obtús}
  \task \rest{pla}
  \task \rest{còncau}
  \task \rest{complet}
  \task \rest{obtús}
  \task \rest{agut}
\end{tasks}

\ex{1} \begin{tasks}(4)
  \task \rest{agut} ($40^\circ$)
  \task \rest{recte} ($90^\circ$)
  \task \rest{obtús} ($140^\circ$)
  \task \rest{còncau} ($230^\circ$)
\end{tasks}

\ex{1} Entre dos nombres de l'esfera hi ha $360^\circ:12=30^\circ$.\begin{tasks}(6)
  \task $\res{90^\circ}$
  \task $\res{120^\circ}$
  \task $\res{150^\circ}$
  \task $\res{180^\circ}$
  \task $\res{90^\circ}$
  \task $\res{30^\circ}$
\end{tasks}

\ex{2} Cada minut gira $360^\circ:60=6^\circ$.\begin{tasks}(4)
  \task $\res{60^\circ}$
  \task $\res{150^\circ}$
  \task $\res{360^\circ}$
  \task $45\cdot6=\res{270^\circ}$
\end{tasks}

\ex{2} \begin{tasks}(1)
  \task $136^\circ:2=\res{68^\circ}$ cadascun.
  \task $2\cdot43^\circ=\res{86^\circ}$ (agut).
\end{tasks}

\ex{2} Entre els dos fan una volta completa: $360^\circ-125^\circ=\res{235^\circ}$.

\ex{3} \begin{tasks}(1)
  \task $360^\circ:5=\res{72^\circ}$; \ $360^\circ:8=\res{45^\circ}$.
  \task $180:15=\res{12}$ vegades.
\end{tasks}


% =====================================================
\seccio{4}{Mesura d'angles: el sistema sexagesimal}

\ex{1} \begin{tasks}(3)
  \task $\res{7200''}$
  \task $\res{900''}$
  \task $\res{4800''}$
  \task $\res{11140''}$
  \task $\res{36610''}$
  \task $\res{2730''}$
\end{tasks}

\ex{1} \begin{tasks}(4)
  \task $\res{300'}$
  \task $\res{150'}$
  \task $\res{765'}$
  \task $\res{10\,800'}$
\end{tasks}

\ex{1} Dividim entre 60 tantes vegades com calga: per exemple, $4000'':60=66'$ i sobren $40''$; $66':60=1^\circ$ i sobren $6'$.\begin{tasks}(3)
  \task $\res{2^\circ\,30'}$
  \task $\res{1^\circ\,6'\,40''}$
  \task $\res{2^\circ\,3'\,5''}$
  \task $\res{2^\circ\,46'\,40''}$
  \task $\res{16^\circ\,40'}$
  \task $\res{24^\circ}$
\end{tasks}

\ex{2} La part decimal dels graus es multiplica per 60: $0,5^\circ=30'$, \ $0,25^\circ=15'$, \ $0,75^\circ=45'$, \ $0,1^\circ=6'$.\begin{tasks}(4)
  \task $\res{12^\circ\,30'}$
  \task $\res{30^\circ\,15'}$
  \task $\res{8^\circ\,45'}$
  \task $\res{6'}$
\end{tasks}

\ex{2} \begin{tasks}(2)
  \task $23^\circ\,45'+18^\circ\,30'=\res{42^\circ\,15'}$
  \task $45^\circ\,36'\,50''+27^\circ\,48'\,25''=\res{73^\circ\,25'\,15''}$
  \task $120^\circ\,15'\,40''+59^\circ\,44'\,20''=\res{180^\circ}$
\end{tasks}

\ex{2} Quan no es pot restar, convertim $1^\circ$ en $60'$ o $1'$ en $60''$.\begin{tasks}(2)
  \task $75^\circ\,40'-32^\circ\,15'=\res{43^\circ\,25'}$
  \task $60^\circ\,20'\,15''-25^\circ\,45'\,30''=\res{34^\circ\,34'\,45''}$
  \task $180^\circ-97^\circ\,34'\,18''=\res{82^\circ\,25'\,42''}$
\end{tasks}

\ex{2} \begin{tasks}(2)
  \task $(12^\circ\,20')\cdot3=\res{37^\circ}$
  \task $(25^\circ\,15'\,40'')\cdot4=\res{101^\circ\,2'\,40''}$
  \task $(8^\circ\,45'\,50'')\cdot6=\res{52^\circ\,35'}$
\end{tasks}

\ex{3} \begin{tasks}(2)
  \task $(75^\circ\,36'):4=\res{18^\circ\,54'}$
  \task $(100^\circ\,20'\,30''):5=\res{20^\circ\,4'\,6''}$
  \task $(47^\circ\,13'):3=\res{15^\circ\,44'\,20''}$
\end{tasks}Comprovació del primer: $(18^\circ\,54')\cdot4=72^\circ\,216'=75^\circ\,36'$ \checkmark

\ex{3} En 1 hora: $360^\circ:24=\res{15^\circ}$. En 1 minut: $15^\circ:60=900':60=\res{15'}$. En 1 segon: $15':60=900'':60=\res{15''}$.


% =====================================================
\seccio{5}{Relacions entre angles}

\ex{1} \begin{tasks}(4)
  \task $\res{60^\circ}$
  \task $\res{43^\circ}$
  \task $\res{1^\circ}$
  \task $\res{74^\circ\,30'}$
\end{tasks}

\ex{1} \begin{tasks}(4)
  \task $\res{135^\circ}$
  \task $\res{60^\circ}$
  \task $\res{1^\circ}$
  \task $\res{87^\circ\,44'\,20''}$
\end{tasks}

\ex{1} $a$ i $c$ són adjacents a l'angle de $35^\circ$: $a=c=180^\circ-35^\circ=\res{145^\circ}$. \ $b$ és oposat pel vèrtex a l'angle de $35^\circ$: $b=\res{35^\circ}$.

\ex{2} L'angle 3 és oposat pel vèrtex a l'1 i els angles 5 i 7 són corresponents a l'1 i al 3: \ $\hat3=\hat5=\hat7=\res{58^\circ}$. \ Els angles 2, 4, 6 i 8 són suplementaris de l'1: \ $180^\circ-58^\circ=\res{122^\circ}$.

\ex{2} \begin{tasks}(1)
  \task \rest{Fals}: dos angles aguts sumen menys de $180^\circ$.
  \task \rest{Cert}: $180^\circ$ menys un angle agut fa més de $90^\circ$.
  \task \rest{Cert}.
  \task \rest{Fals}: ja passa de $90^\circ$.
\end{tasks}

\ex{2} $89^\circ\,59'\,60''-37^\circ\,25'\,48''=\res{52^\circ\,34'\,12''}$ \qquad $179^\circ\,59'\,60''-37^\circ\,25'\,48''=\res{142^\circ\,34'\,12''}$

\ex{3} $x+(x+20)=90\ \Rightarrow\ 2x=70\ \Rightarrow\ x=35$. \ Els angles mesuren \rest{$35^\circ$ i $55^\circ$}.

\ex{3} $x+4x=180\ \Rightarrow\ 5x=180\ \Rightarrow\ x=\res{36^\circ}$ (el suplementari fa $144^\circ=4\cdot36^\circ$).

\ex{3} Els angles alterns interns són iguals: $3x+10=70\ \Rightarrow\ 3x=60\ \Rightarrow\ x=\res{20^\circ}$.


% =====================================================
\seccio{6}{Polígons: elements i classificació}

\ex{1} \begin{tasks}(3)
  \task \rest{pentàgon}
  \task \rest{octàgon}
  \task \rest{hexàgon}
  \task \rest{decàgon}
  \task \rest{heptàgon}
  \task \rest{dodecàgon}
\end{tasks}

\ex{1} \begin{tasks}(2)
  \task 5 costats, \rest{convex} i \rest{regular}
  \task 5 costats, \rest{còncau} i \rest{irregular}
  \task 4 costats, \rest{convex} i \rest{irregular} (rectangle)
  \task 4 costats, \rest{convex} i \rest{irregular}
\end{tasks}

\ex{1} De cada vèrtex ixen $n-3$ diagonals; en total, $\dfrac{n\cdot(n-3)}{2}$.\begin{tasks}(4)
  \task 1 i \res{2}
  \task 2 i \res{5}
  \task 5 i \res{20}
  \task 7 i \res{35}
\end{tasks}

\ex{2} \begin{tasks}(1)
  \task \rest{Cert} (i també tots els angles iguals).
  \task \rest{Fals}: els seus angles no són tots iguals.
  \task \rest{Cert}.
  \task \rest{Fals}: no en té cap ($\frac{3\cdot0}{2}=0$).
\end{tasks}

\ex{2} Busquem $n$ amb $n\cdot(n-3)=28$: \ $7\cdot4=28$ \ $\Rightarrow$ \rest{heptàgon}. \ Amb $n\cdot(n-3)=54$: \ $9\cdot6=54$ \ $\Rightarrow$ \rest{enneàgon}.

\ex{3} $n-3=9\ \Rightarrow\ n=\res{12}$ (dodecàgon). \ En total: $\dfrac{12\cdot9}{2}=\res{54}$ diagonals.


% =====================================================
\seccio{7}{Triangles}

\ex{1} \begin{tasks}(4)
  \task \rest{equilàter}
  \task \rest{isòsceles}
  \task \rest{escalé}
  \task \rest{isòsceles}
\end{tasks}

\ex{1} \begin{tasks}(2)
  \task $180^\circ-100^\circ=\res{80^\circ}$, \rest{acutangle}
  \task $180^\circ-125^\circ=\res{55^\circ}$, \rest{rectangle}
  \task $180^\circ-145^\circ=\res{35^\circ}$, \rest{obtusangle}
  \task $180^\circ-90^\circ=\res{90^\circ}$, \rest{rectangle} (i isòsceles)
\end{tasks}

\ex{1} El costat més llarg ha de ser més menut que la suma dels altres dos.\begin{tasks}(4)
  \task $5<7$: \rest{sí}
  \task $6>5$: \rest{no}
  \task $9<10$: \rest{sí}
  \task $8=8$: \rest{no}
\end{tasks}

\ex{2} \begin{tasks}(1)
  \task $(180^\circ-50^\circ):2=\res{65^\circ}$ cadascun.
  \task $180^\circ-2\cdot72^\circ=\res{36^\circ}$.
  \task $(180^\circ-90^\circ):2=\res{45^\circ}$ cadascun.
\end{tasks}

\ex{2} \begin{tasks}(2)
  \task \rest{Sí}: angles de $90^\circ$, $45^\circ$ i $45^\circ$.
  \task \rest{No}: els angles d'un equilàter fan $60^\circ$.
  \task \rest{Sí}: per exemple, $50^\circ$, $60^\circ$ i $70^\circ$.
  \task \rest{No}: els angles d'un equilàter fan $60^\circ$.
\end{tasks}

\ex{2} Els dos angles aguts sumen $90^\circ$: \ $89^\circ\,60'-27^\circ\,35'=\res{62^\circ\,25'}$.

\ex{2} \begin{tasks}(1)
  \task El \rest{circumcentre} (tall de les mediatrius).
  \task L'\rest{ortocentre}.
  \task El \rest{baricentre} (tall de les mitjanes).
  \task L'\rest{incentre} (tall de les bisectrius).
\end{tasks}

\ex{3} $x+(x+15)+(x+30)=180\ \Rightarrow\ 3x+45=180\ \Rightarrow\ x=45$. \ Angles: \rest{$45^\circ$, $60^\circ$ i $75^\circ$}: \rest{acutangle} i escalé.

\ex{3} El tercer costat ha de ser més menut que $5+12=17$ i més gran que $12-5=7$: \ \rest{8, 9, 10, …, 16\,cm} (9 possibilitats).

\ex{3} L'angle adjacent fa $180^\circ-130^\circ=50^\circ$. \ El tercer: $180^\circ-60^\circ-50^\circ=70^\circ$. \ Angles: \rest{$50^\circ$, $60^\circ$ i $70^\circ$}.


% =====================================================
\seccio{8}{Quadrilàters}

\ex{1} \begin{tasks}(3)
  \task \rest{quadrat}
  \task \rest{rombe}
  \task \rest{trapezi isòsceles}
  \task \rest{rectangle}
  \task \rest{romboide}
  \task \rest{trapezi rectangle}
\end{tasks}

\ex{1} \begin{tasks}(3)
  \task $\res{80^\circ}$
  \task $\res{105^\circ}$
  \task $\res{85^\circ}$
\end{tasks}

\ex{1} \begin{tasks}(2)
  \task \rest{quadrat i rectangle}
  \task \rest{quadrat i rombe}
  \task \rest{quadrat i rectangle}
  \task \rest{quadrat i rombe}
\end{tasks}

\ex{2} \begin{tasks}(1)
  \task $\res{54^\circ,\ 126^\circ,\ 54^\circ,\ 126^\circ}$ (angles oposats iguals; consecutius suplementaris)
  \task $\res{112^\circ,\ 68^\circ,\ 112^\circ,\ 68^\circ}$
  \task $\res{90^\circ,\ 90^\circ,\ 63^\circ,\ 117^\circ}$ ($180^\circ-63^\circ=117^\circ$)
\end{tasks}

\ex{2} \begin{tasks}(1)
  \task \rest{Cert}: té els quatre angles rectes.
  \task \rest{Fals}: el rombe no té per què tindre angles rectes.
  \task \rest{Cert}.
  \task \rest{Fals}: si en té tres, el quart també fa $90^\circ$ i seria un rectangle.
\end{tasks}

\ex{3} \begin{tasks}(1)
  \task $10x=360\ \Rightarrow\ x=36$: \ \rest{$36^\circ$, $72^\circ$, $108^\circ$ i $144^\circ$}.
  \task Dos angles consecutius sumen $180^\circ$: $x+3x=180\Rightarrow x=45$. \ \rest{$45^\circ$, $135^\circ$, $45^\circ$ i $135^\circ$}.
\end{tasks}


% =====================================================
\seccio{9}{Angles dels polígons. Polígons regulars}

\ex{1} \begin{tasks}(5)
  \task $\res{540^\circ}$
  \task $\res{720^\circ}$
  \task $\res{1080^\circ}$
  \task $\res{1440^\circ}$
  \task $\res{1800^\circ}$
\end{tasks}

\ex{1} \begin{tasks}(3)
  \task $\res{60^\circ}$ i $\res{120^\circ}$
  \task $\res{90^\circ}$ i $\res{90^\circ}$
  \task $\res{108^\circ}$ i $\res{72^\circ}$
  \task $\res{120^\circ}$ i $\res{60^\circ}$
  \task $\res{135^\circ}$ i $\res{45^\circ}$
  \task $\res{144^\circ}$ i $\res{36^\circ}$
\end{tasks}

\ex{2} $n=360:24=\res{15}$ costats. \ Angle interior: $180^\circ-24^\circ=\res{156^\circ}$ \ (o bé $\frac{13\cdot180^\circ}{15}=156^\circ$).

\ex{2} Angle central: $180^\circ-160^\circ=20^\circ$; \ $n=360:20=\res{18}$ costats.

\ex{2} $(n-2)\cdot180=2340\ \Rightarrow\ n-2=13\ \Rightarrow\ n=\res{15}$ costats.

\ex{3} Amb $100^\circ$, l'angle central faria $80^\circ$ i $360:80=4,5$ no és un nombre natural: \rest{no}. \ Amb $150^\circ$: $360:30=12$: \rest{sí}, el dodecàgon regular.

\ex{3} Els angles que es troben en un vèrtex han de sumar $360^\circ$: \rest{triangles} ($6\cdot60^\circ$), \rest{quadrats} ($4\cdot90^\circ$) i \rest{hexàgons} ($3\cdot120^\circ$). Amb pentàgons no ($3\cdot108^\circ=324^\circ$) ni amb octàgons sols. Combinant-los, \rest{sí}: dos octàgons i un quadrat, $135^\circ+135^\circ+90^\circ=360^\circ$.

\ex{3} Cada angle central fa $360^\circ:6=60^\circ$. Els dos radis són iguals, així que els altres dos angles són iguals: $(180^\circ-60^\circ):2=60^\circ$. Els sis triangles són \rest{equilàters}.


\end{document}
